\section{Results}
\label{sec:results}

We present results by research question. All numbers derive from the 33 S2AG-resolved papers from the huashen218 proxy corpus.

\subsection{R1: Pipeline Infrastructure Validated (23/23 Tests Pass)}

The complete S2AG bibliometric pipeline passes all 23 unit tests in 0.85 seconds. Zero failures.
The test suite covers corpus ingestion, S2AG resolution, classification (all three schemes), graph construction, gate evaluation, and figure generation.

\textbf{What this means:} The gate failure reported below is a data access failure, not a pipeline failure.
Every algorithmic component is correct and validated. The infrastructure contribution is complete.

\subsection{R2: FoS-Primary Classification Achieves 100\% Coverage vs.\ 12\% for Venue Strings}

This is the central methodological result.

\begin{table}[h]
\centering
\caption{Classification coverage by scheme.}
\label{tab:coverage}
\begin{tabular}{llcccc}
\toprule
Scheme & Type & Classified & Coverage & ML\_NLP & HCI \\
\midrule
Scheme 1 & Venue string (broad) & 4/33 & 12.1\% & 3 & 1 \\
Scheme 2 & Venue string (narrow) & 4/33 & 12.1\% & 3 & 1 \\
\textbf{Scheme 3} & \textbf{FoS-primary} & \textbf{33/33} & \textbf{100.0\%} & \textbf{28} & \textbf{5} \\
\bottomrule
\end{tabular}
\end{table}

Schemes 1 and 2 produce identical coverage because all classifiable papers have abbreviated venue names in the target sets --- and the 29 unclassified papers all have full proceedings names in S2AG that do not match any substring.

\textbf{Why this gap exists:} S2AG's \texttt{venue} field stores the full proceedings name as submitted by the publisher.
Substring matching against ``NeurIPS'' fails silently on ``Proceedings of the 37th Annual Conference on Neural Information Processing Systems.''
The \texttt{fieldsOfStudy} field is format-independent and consistently populated.
The 12\% coverage rate is a format mismatch artifact, not a corpus composition artifact.

\textbf{Implication:} Any bibliometric study of an interdisciplinary corpus using S2AG data and venue-string classification will silently drop the majority of papers unless FoS-primary classification is used.

\subsection{R3: Directional Citation Signal --- Ratio $\approx$ 0.107, Gate Thresholds Not Met}

\textbf{Coverage and Gate Evaluation:}

\begin{itemize}
\item S2AG resolution rate: 33/49 = \textbf{67.3\%} (gate: 70\%) --- gate FAIL by 2.7 pp
\item Cross-group edge counts per scheme:
\end{itemize}

\begin{table}[h]
\centering
\caption{Cross-group directed edge counts per scheme.}
\label{tab:edges}
\begin{tabular}{lccc}
\toprule
Scheme & ML\_NLP$\to$HCI & HCI$\to$ML\_NLP & Total Cross-group \\
\midrule
Scheme 1 & 1 & 1 & 2 \\
Scheme 2 & 1 & 1 & 2 \\
Scheme 3 & 5 & 4 & \textbf{9} \\
\bottomrule
\end{tabular}
\end{table}

Maximum cross-group edges: \textbf{9} (gate: 30) --- gate FAIL by 21 edges.

Both gate thresholds are not satisfied.
Statistical tests (P1: chi-squared, P2: proportion z-test, P3: betweenness centrality) are not executed.
All three predictions are INCONCLUSIVE --- not REFUTED.
Figure~\ref{fig:gate} shows the gate evaluation results across all three schemes.

\textbf{Directed Edge Matrix (Scheme 3):}

\begin{table}[h]
\centering
\caption{Directed citation edge matrix under Scheme 3 (FoS-primary).}
\label{tab:matrix}
\begin{tabular}{lcc c}
\toprule
 & $\to$ ML\_NLP & $\to$ HCI & Outgoing total \\
\midrule
\textbf{ML\_NLP $\to$} & 42 & 5 & 47 \\
\textbf{HCI $\to$} & 4 & 0 & 4 \\
\bottomrule
\end{tabular}
\end{table}

Citation proportions: ML\_NLP$\to$HCI = 5/47 = \textbf{10.6\%}; HCI$\to$ML\_NLP = 4/4 = \textbf{100.0\%}; ratio = \textbf{0.107}.

Figure~\ref{fig:heatmaps} shows the directed edge count heatmaps for all three classification schemes.
The Scheme 3 panel is the primary result.

\textbf{Interpretation:} The directional signal is directionally consistent with H-CitAsym.
ML/NLP alignment papers allocate 42 of 47 within-corpus outgoing citations (89.4\%) to other ML/NLP papers, with only 5 citations (10.6\%) crossing to HCI.
HCI alignment papers allocate 100\% of within-corpus outgoing citations to ML/NLP --- no HCI-to-HCI within-corpus citations appear.
The ratio of 0.107 is far below the null hypothesis value of 1.0.
No falsifier was triggered.
\textbf{Important caveat:} At $N=9$ cross-group edges (5 from ML/NLP papers, 4 from HCI papers with within-corpus citations), this ratio is not statistically confirmable.
Of the 5 HCI papers resolved under Scheme 3, those with outgoing within-corpus citations all reference ML/NLP work exclusively.

\textbf{Unexpected Finding: HCI$\to$HCI = 0.}
The complete absence of HCI-to-HCI within-corpus citations was not anticipated even under H-CitAsym.
At $N=5$ HCI papers this could be a sample artifact, but the pattern suggests that HCI alignment papers treat ML/NLP alignment work as their primary reference frame, not prior HCI alignment work.

\subsection{Corpus Proxy Analysis}

Figure~\ref{fig:dropout} shows unresolved papers by ID type.
All 16 unresolved papers are ACM Digital Library publications without arXiv preprints.
The dropout is non-systematic: both ML/NLP and HCI papers appear in the unresolved set --- no directional bias is introduced.

The corpus extraction funnel: $\sim$130 linked papers $\to$ 49 extractable identifiers (38\% extraction rate) $\to$ 33 S2AG-resolved papers (67.3\% resolution rate).
Figure~\ref{fig:pie} shows corpus composition under Scheme 1.

\begin{table}[h]
\centering
\caption{Summary of results.}
\label{tab:summary}
\begin{tabular}{lll}
\toprule
Research Question & Result & Confidence \\
\midrule
RQ1: Pipeline validated? & 23/23 tests pass & HIGH \\
RQ2: FoS-primary coverage? & 100\% vs.\ 12\% (Scheme 1/2) & HIGH \\
RQ3: Gate thresholds met? & No (coverage=67.3\%, edges=9) & HIGH \\
RQ3: Directional ratio? & 0.107 (consistent with H-CitAsym) & MEDIUM ($N=9$) \\
P1/P2/P3 statistical tests? & INCONCLUSIVE (gate not met) & --- \\
\bottomrule
\end{tabular}
\end{table}
